\section{Introduction}
\label{sec:introduction}

Consider two integral affine forms
\begin{equation}
 L_1(r)=d+sr,\qquad
 L_2(r)=u+ar
\label{eq:affine-pair}
\end{equation}
whose coefficient determinant is prescribed:
\begin{equation}
 su-ad=h_0\ne0.
\label{eq:fixed-determinant}
\end{equation}
This is the literal algebraic shape returned by the primitive-column
gate of the fixed-shift TPC route \citep{WangTPC33}.  The same
determinant appears in the two-point logarithmic Chowla setting, but
the physical packet has growing coefficients, ordinary weighted
fibers, prescribed origins, masks, and a nontrivial outer
reassembly.  Those quantifiers cannot be removed by notation.

The present paper asks a deliberately narrower question:
\begin{quote}
\emph{Which part of the two-Mobius coefficient is completely
explained by primitive affine local arithmetic, and which part
remains a genuine parity correlation?}
\end{quote}
The answer is exact.  The fixed determinant confines every common
value divisor to \(h_0\).  Primitivity then gives a closed formula for
the root count modulo every prime power.  At the squarefree level,
the forbidden residues form an explicit periodic local support with
a positive Euler product.  None of these facts cancels the remaining
Liouville signs.

\subsection{The three proof levels}

We use the proof-level convention of the surrounding TPC series.
\begin{description}[leftmargin=3.2em,style=nextline]
\item[\(\Lzero\):]
finite integer algebra, exact congruence counts, Chinese-remainder
identities, fixed-form density theorems, and structure-free logical
separations;
\item[\(\Lone\):]
verification that the exact objects are the literal physical
coefficients on their actual intervals, with every mask, origin,
weight, polarization, key, and global normalization retained;
\item[\(\Ltwo\):]
a new uniform arithmetic cancellation estimate for the resulting
growing fixed-\(h_0\) family.
\end{description}
This paper proves \(\Lzero\) theorems and formulates an exact
\(\Lone\) crosswalk.  It proves no \(\Ltwo\) estimate.

\subsection{Main results}

Assume first that the forms are primitive:
\[
 \gcd(d,s)=\gcd(u,a)=1.
\]
For a prime \(p\) and \(k\ge1\), write
\[
 \Rset_{p,k}
 =
 \{r\bmod p^k:p^k\mid L_1(r)\}
 \cup
 \{r\bmod p^k:p^k\mid L_2(r)\},
 \qquad
 \nu_{p,k}=\#\Rset_{p,k}.
\]
The exact count is
\begin{equation}
 \nu_{p,k}
 =
 \ind_{p\nmid s}+\ind_{p\nmid a}
 -\ind_{p\nmid as}\ind_{p^k\mid h_0}.
\label{eq:intro-root-count}
\end{equation}
On the two-nonconstant branch \(as\ne0\), this formula distinguishes
two kinds of exceptional prime:
\begin{enumerate}[label=(\roman*),leftmargin=2.2em]
\item a \emph{joint} or collision prime divides \(h_0\), hence belongs
to a fixed finite set when \(h_0\) is fixed;
\item a one-slope prime divides exactly one of \(a,s\), may move with
the growing coefficients, and deletes one local root rather than
creating a simultaneous obstruction.
\end{enumerate}
In particular, a collision modulo \(p\) requires \(p\mid h_0\), while
a collision modulo \(p^2\) requires the strictly stronger condition
\(p^2\mid h_0\).

The squarefree support is
\[
 \Qsf(r)=\mu^2(|L_1(r)|)\mu^2(|L_2(r)|),
\]
away from the at most two zeros of the forms.  Its local density is
\[
 \Mloc(L_1,L_2)
 =
 \prod_p\left(1-\frac{\nu_{p,2}}{p^2}\right)>0.
\]
For fixed coefficients this is the natural density of
\(\Qsf(r)=1\).  This is a special linear instance of the classical
squarefree-pattern framework of \citet{Mirsky1949}; we include a
complete elementary proof because the exact dependence on the
slopes and \(h_0\) matters here.

Finally, on a positive active interval,
\[
 \mu(L_1)\mu(L_2)
 =
 \Qsf\,\lambda(L_1)\lambda(L_2).
\]
The local factor \(\Qsf\) is fully classified.  The surviving
weighted correlation of the two Liouville signs is not.  Tao's
two-point theorem treats fixed affine forms with logarithmic
averaging \citep{Tao2016}; averaged-shift theorems change a different
quantifier \citep{MatomakiRadziwillTao2015}.  Neither theorem supplies
the literal growing-coefficient, fixed-\(h_0\), ordinary weighted,
uniform-prefix statement required here; the logarithmic-to-ordinary
gap is audited in TPC-87 \citep{WangTPC87}.

\subsection{What the split does not say}

A positive squarefree Euler product says that no finite collection
of square-power congruences annihilates the support.  It does not
give cancellation of \(\lambda(L_1)\lambda(L_2)\), does not select a
prescribed shift from an averaged theorem, and does not produce a
lower bound for primes.  In a growing physical family, even the
passage from a formal local product to uniform support mass on every
short translated fiber requires its own estimate.

The conclusion is therefore a precise fork.  A positive continuation
must prove sign cancellation on the literal squarefree support.  A
negative continuation may show that a proposed theorem cannot retain
the growing forms, interval origins, or endpoint budget.  Either
outcome must keep the compatibility and physical-loss ledgers
independent.

\subsection{Organization}

\Cref{sec:determinant-algebra} proves the determinant and content
identities.  \Cref{sec:local-root-counts} gives the complete
prime-power root classification.  The squarefree Euler product,
fixed-form density, and uniformity warning appear in
\cref{sec:squarefree-support}.  \Cref{sec:parity-split} proves the
exact Mobius--Liouville and square-divisor decompositions.
\Cref{sec:physical-interface} identifies the unresolved literal
correlation.  The independent ledgers and final claim boundary are
recorded in \cref{sec:separate-ledgers,sec:claim-boundary}.
